% Drawing 80/02-8 -- folder 80, batch 02, archive page 32 (PDF page 33).
% Transcription: à gauche, un dessin : trois droites en position générale formant un triangle
% Kind: three lines in general position. Method 2 (schematic redraw from the facsimile). Status: ok.
% Notes: Ink. Three slightly curved lines through their three computed crossing points. Redrawn 2026-11-12 (issue 128): to check again.
% Coordinates: 1 unit = 1 mm = 1/100 of the crop width (bbox in page fractions: [0.12, 0.47, 0.42, 0.63]).
% Redrawn by Claude (claude-fable-5-2), issue #128, 2026-11-12; not checked by a human; not compiled (no TeX in the session).
% Rights: reproduces part of the work -- behind RENDER_TRANSCRIPTS, never published by itself.
\documentclass[tikz,border=3pt]{standalone}
\usepackage{amsmath,amssymb}
\usetikzlibrary{patterns,arrows.meta}
\tikzset{pencil/.style={black!45}, struck/.style={black!35}, mpt/.style={circle,draw,inner sep=1.3pt}, pt/.style={circle,fill,inner sep=0.9pt}, opt/.style={circle,draw,fill=white,inner sep=1.1pt}}
\begin{document}
\begin{tikzpicture}[x=1mm,y=1mm,line cap=round,line join=round,font=\scriptsize]
\draw plot[smooth, tension=0.55] coordinates {(24,-4.5) (30,-11) (40,-20) (48.5,-27.5) (61,-38) (75,-49) (85,-58) (92,-66)};
\draw plot[smooth, tension=0.55] coordinates {(57.5,-13.5) (53,-21) (48.5,-27.5) (39,-39.3) (30,-50) (21,-59.5) (12,-70)};
\draw plot[smooth, tension=0.55] coordinates {(5.5,-41) (15,-41.3) (30,-40.3) (39,-39.3) (50,-38.5) (61,-38) (75,-36.3) (89,-34.8)};
\end{tikzpicture}
\end{document}
